% Pursuing Stacks (cotes 134-2 to 134-8): archive cote 134-8, site page 32 (PDF page 33), typescript page 586.
% typescript. Draft compared: ps7.tex:1496-1554.
% AI-transcribed by Claude Opus 5.5 (claude-opus-5-5, adaptive thinking, effort high), Message Batches API, batch ps4 (msgbatch_012Aoc1NuskATPsyBeKSYK7J), 2026-11-12; not yet checked by eye.
\scanpage{134-8/32}
\mspage{586}

thus the Modules
(12) \qquad $\underline{O}_A^{\;(a)}$ \quad , for a in A \quad ,

are clearly projective, and \unsure{there|theire} are ``sufficiently many''. Thus,
any \inserted{additive} functor from Mod($\underline{O}_A$) to an abelian category admits a total left
derived functor. However, it is not always true, for a map (1)
of small ringed categories, that the corresponding inverse image
functor for Modules
\[ G \longmapsto f^{*}(G)= f_o^{\;*}(G) \otimes_{f_o^{*}(\underline{O}_B)} \underline{O}_A \]
admits \xout{always} a left adjoint $f_{!}$ , or what amounts to the same\xout{x}, that
this functor (which is right exact) is left exact and commutes to
\xout{arbitrary} \inserted[typed]{small} products. It is'nt even true necessarily, when we assume
A and B to be the final category ! Left exactness of $f^{\hand{*}}$ just means
flatness of $\underline{O}_A$ over $\underline{O}_B$ , i.e. that for any a in A, $\underline{O}_A(a)$ is flat as
a\xout{s} module over $\underline{O}_B(b)$, where b=f(a). As for commutation to small
products, it amounts (\xout{\unsure{when}} \inserted{together with} the first condition)\xout{is satisfied)} to the
still more ex\hand{a}cting condition that $\underline{O}_A(a)$ should be a \uline{projective module
of finite type over} $\underline{O}_B(b)$, for any a in A. This condition is so close
to the condition $\underline{O}_{\hand{A}} = f_o^{\;*}(\underline{O}_B)$ \quad (already considered in (11) above),
that for a bird's eye view as we are aiming at here, we may as well
assume this slightly stronger condition\hand{!} Anyhow, as noticed before,
to get \xout{compatibility of} commutation of $f_{!}$ with restriction of scalars
and with colocalization, even this assumption is'nt enough apparently\hand{,}
and we'll have to assume moreover that $\underline{O}_B$ (hence also $\underline{O}_A$) is constant,
or for the least, locally constant.

This brings us back to the case when a fixed ring k is given,
and when we are working with categories of presheaves of k-modules
- a situation studied at length in chapter V. In case the target
category B is the final category, hence L$f_{!}$ is just (absolute) total
homology of A, this then may be computed nicely, using an ``integra-%
tor'' for \hand{A} , namely a projective resolution of $k_{A^o}$ in $(A^o)^{\wedge}_{k}$ . This
is more or less where we ended up \inserted{with} by the end of chapter V, when \over{d}{d}eve-%
\xout{\unsure{were}} lopping a nicely autodual homology-cohomology set-up, \xout{where}
replacing the category of k-modules $\mathrm{Ab}_k$ by a more or less arbitrary
abelian category. I was about to \xout{\unsure{carry thr}} go on and carry th\hand{r}ough
a similar treatment in the ``r\over{\unsure{ee}}{el}ative'' case, namely for an arbitrary
map f in (Cat) \xout{\unsure{,}}\hand{(}but then I got caught unsuspectingly by that unen-%
ding digression on schematic homotopy types, finally making up a whole
chapter by itself\hand{)}. Maybe it is still worthwhile to come back to this\hand{,}

\margin{32}
